\chapter{Funkcionálne požiadavky}

\section{Obrazovky aplikácie}

Aplikácia obsahuje nasledujúcich 7 funkčných obrazoviek (splash screen a login screen sa do počtu nepočítajú):

\begin{enumerate}
  \item \textbf{Job List} — zoznam jobov (Manager: všetky; Worker: pridelené)
  \item \textbf{Job Detail} — detail jobu, stav, popis, fotky, tlačidlo na zmenu stavu
  \item \textbf{Create / Edit Job} — formulár na vytvorenie alebo úpravu jobu (len Manager)
  \item \textbf{Map View} — mapa s polohami pracovníkov a pinmi jobov
  \item \textbf{Chat} — real-time konverzácia (per job alebo globálna)
  \item \textbf{Worker List} — zoznam pracovníkov a ich aktuálny stav (len Manager)
  \item \textbf{Profile / Settings} — profil používateľa, odhlásenie, prepínanie dark mode
\end{enumerate}

\section{CRUD operácie}

\begin{table}[H]
  \centering
  \begin{tabularx}{\textwidth}{llX}
    \toprule
    \textbf{Entita} & \textbf{Operácia} & \textbf{Popis} \\
    \midrule
    Job & Create  & Manažér vytvorí nový job cez formulár \\
    Job & Read    & Zoznam a detail jobov pre obe roly \\
    Job & Update  & Manažér upraví job; Worker zmení stav \\
    Job & Delete  & Manažér zmaže job \\
    \bottomrule
  \end{tabularx}
  \caption{CRUD operácie nad entitou Job}
\end{table}

\section{Offline režim}

Aplikácia podporuje offline fungovanie pomocou lokálnej SQLite databázy (\textbf{WatermelonDB}). Keď zariadenie nemá internetové pripojenie:

\begin{itemize}
  \item Worker môže zmeniť stav jobu (Create, Update, Delete stavu/poznámky) — zmena sa uloží lokálne s príznakom \texttt{pending\_sync}.
  \item Po obnovení pripojenia sa zmeny automaticky synchronizujú so serverom na pozadí.
  \item Konflikt (napr. job bol medzitým zmazaný manažérom) sa rieši politikou \textit{server wins}: server vráti chybu, aplikácia zobrazí používateľovi upozornenie a lokálnu zmenu zahodí.
\end{itemize}

\section{WebSockets — real-time komunikácia}

Aplikácia využíva WebSocket spojenie (Socket.IO) pre:

\begin{itemize}
  \item \textbf{Live chat} — obojsmerné zasielanie správ medzi Manažérom a Workerom v rámci jobu.
  \item \textbf{Zdieľanie polohy} — Worker periodicky (každých 10 s) odosiela svoju GPS polohu; Manažér ju dostáva v reálnom čase a vidí ju na mape.
\end{itemize}

\section{Push notifikácie}

\begin{table}[H]
  \centering
  \begin{tabularx}{\textwidth}{llX}
    \toprule
    \textbf{Trigger} & \textbf{Príjemca} & \textbf{Správanie} \\
    \midrule
    Nový job pridelený (server) & Worker & Notifikácia na pozadí; klik presmeruje na Job Detail \\
    Job dokončený (app)         & Manager & Notifikácia na pozadí; klik presmeruje na Job Detail \\
    \bottomrule
  \end{tabularx}
  \caption{Push notifikácie}
\end{table}

\section{Permission handling}

Aplikácia si vyžiada nasledujúce systémové oprávnenia:

\begin{itemize}
  \item \textbf{Kamera} — fotografovanie dokumentácie k jobu
  \item \textbf{Galéria / Súbory} — výber existujúcej fotografie
  \item \textbf{Poloha (foreground + background)} — zdieľanie GPS polohy Workera
  \item \textbf{Push notifikácie} — prijímanie FCM správ
\end{itemize}

\section{Ďalšie povinné požiadavky}

\begin{itemize}
  \item \textbf{Dark mode} — aplikácia plne podporuje systémový dark mode aj manuálny prepínač v Settings.
  \item \textbf{User management} — perzistentná autentifikácia pomocou JWT tokenu uloženého v bezpečnom úložisku (\texttt{react-native-keychain}).
\end{itemize}